% ===========================================================================
% justificacion.tex  --  Capitulo 3
% Se incluye desde main.tex con \input{capitulos/justificacion}
% ===========================================================================
\section{Justificación}
\label{cap:justificacion}

\subsection{Relevancia del problema de salud}
\label{sec:just-salud}

La desregulación del sistema serotoninérgico se asocia con un conjunto amplio de
trastornos psiquiátricos y neurológicos: migraña, dolor, ansiedad,
esquizofrenia, trastornos del movimiento y depresión; esta última afecta a más
de 300 millones de personas en el mundo, sin distinción de género, origen o
condición socioeconómica, y en su forma avanzada puede conducir al deterioro
grave de la salud y de la calidad de vida \cite{czub_curated_2021}.

El receptor 5-HT\textsubscript{1A} ocupa una posición privilegiada dentro de ese
sistema: es el subtipo mejor caracterizado de los catorce descritos y participa
en la modulación de la transmisión neuronal, la plasticidad sináptica y la
neurogénesis \cite{czub_curated_2021}. Varios fármacos en uso clínico actúan sobre
él aripiprazol, buspirona, clozapina, risperidona, vortioxetina, entre
otros, pero los medicamentos disponibles presentan numerosos efectos
adversos, lo que mantiene abierta la búsqueda de estructuras con mejor perfil
terapéutico \cite{czub_automlbased_2022}. Existe además evidencia de que el receptor
podría ser blanco en terapia oncológica, lo que amplía el valor de disponer de
modelos predictivos de afinidad para esta diana \cite{czub_curated_2021}.

\subsection{Por qué un enfoque \textit{in silico}}
\label{sec:just-insilico}

Los métodos experimentales de referencia son lentos y caros. En el caso concreto
de la permeabilidad de la barrera hematoencefálica, el \textit{log\,BB} en
roedor es uno de los estándares \textit{in vivo}, y como la mayoría de los
métodos \textit{in vivo} consume tiempo, recursos y animales de experimentación
\cite{faramarzi_development_2022}. Los modelos predictivos basados en estructura permiten
priorizar qué compuestos merecen medición experimental, reduciendo tanto el
costo como el uso de animales \cite{faramarzi_development_2022,spielvogel_enhancing_2025a}.

El argumento no es que el modelo sustituya al experimento. No lo hace, y
sostener lo contrario sería insostenible ante cualquier revisión seria. El
argumento es de \textbf{asignación de recursos}: ante un espacio de búsqueda de
hasta $10^{60}$ moléculas con propiedades de fármaco y bibliotecas virtuales que
ya superan los 36 mil millones de compuestos \cite{hasselgren_artificial_2024}, el cuello
de botella no es la capacidad de ensayar, sino la capacidad de decidir qué
ensayar. Un modelo que ordena candidatos mejor que el azar tiene valor aunque su
exactitud absoluta sea moderada.

\subsection{Por qué integrar tres propiedades y no solo la afinidad}
\label{sec:just-integracion}

Ésta es la aportación central del trabajo, y descansa sobre un hecho empírico.
La distribución de causas de fracaso muestra que la eficacia y la seguridad
fallan en proporciones comparables: en Fase II, entre el 50\,\% y el 60\,\% de
los fracasos corresponden a falta de eficacia y alrededor del 30\,\% a
problemas de seguridad \cite{mao_artificial_2026}; antes de la clínica, cerca del
40\,\% de los candidatos preclínicos se descarta por perfiles ADMET
insuficientes \cite{zhang_computational_2025}.

El caso DSP-1181, es la ilustración más directa: un agonista del receptor 5-HT\textsubscript{1A} diseñado con
inteligencia artificial alcanzó la etapa \textit{hit-to-lead} en menos de doce
meses y fue descontinuado en Fase I por prolongación del intervalo QT detectada
en estudios preclínicos de seguridad \cite{mao_artificial_2026}. Es decir, un programa
que optimizó afinidad con enorme eficiencia fracasó por una propiedad que no
estaba en su función objetivo. Un análisis sistemático de veinte programas de
descubrimiento asistido por IA entre 2018 y 2023 reporta el mismo patrón: la
optimización asistida redujo el tiempo de cribado entre 30\,\% y 50\,\%, pero la
tasa de abandono en Fases I y II se mantuvo comparable a la del resto de la
industria, con la seguridad y las propiedades farmacocinéticas inesperadas como
causas principales \cite{mao_artificial_2026}.

De aquí se sigue la justificación del diseño propuesto: evaluar afinidad,
permeabilidad de la BHE y toxicidad en un mismo flujo, sobre una representación
molecular común, para que las tres restricciones actúen simultáneamente durante
la priorización y no de manera secuencial. Se trata de tres propiedades que
\textbf{condicionan} la viabilidad de un compuesto como fármaco del sistema
nervioso central; ninguna de ellas la determina por sí sola.

\subsection{Por qué redes neuronales de grafos, y por qué con línea base}
\label{sec:just-gnn}

Las moléculas admiten una representación natural como grafos, con átomos como
nodos y enlaces como aristas, y las redes neuronales de grafos están construidas
precisamente para operar sobre esa estructura \cite{mao_artificial_2026}. La ventaja
conceptual es que la representación se aprende junto con la tarea en lugar de
fijarse de antemano, y que la misma familia de arquitecturas admite tanto
regresión de afinidad como clasificación multitarea de propiedades ADMET, lo
que hace viable un \textit{pipeline} con una representación compartida.

Ahora bien, la elección del método no debe confundirse con una apuesta. Los
enfoques clásicos de QSAR, basados en descriptores documentados y algoritmos
transparentes, siguen produciendo modelos competitivos y auditables: un modelo
de AutoML sobre descriptores Mordred, entrenado con cerca de 9\,500 ligandos del
receptor 5-HT\textsubscript{1A}, alcanza desempeño de referencia y satisface
cuatro de los cinco principios OECD de validación \cite{czub_automlbased_2022}. En
predicción de permeabilidad de la BHE, un bosque aleatorio sobre parámetros
moleculares supera con claridad a las reglas heurísticas establecidas
\cite{spielvogel_enhancing_2025a}. En conjuntos de tamaño moderado, una línea base bien
ajustada frecuentemente iguala o supera a una arquitectura profunda.

Por ello, el trabajo incorpora la comparación contra líneas base fuertes como
requisito metodológico permanente y no como etapa preliminar. Si la línea base
gana, ése es un resultado legítimo del estudio y se reporta como tal. Un trabajo
que solo reporta el caso en que su método favorito gana no es un estudio; es una
demostración.

\subsection{Por qué el protocolo de evaluación es parte de la aportación}
\label{sec:just-evaluacion}

Buena parte del valor —y de la fragilidad— de estos modelos reside en cómo se
evalúan, no en qué arquitectura usan.

En primer lugar, la partición de los datos. Si compuestos estructuralmente
análogos quedan repartidos entre entrenamiento y prueba, el modelo resuelve un
problema de interpolación y las métricas obtenidas no anticipan su desempeño
sobre quimiotipos nuevos, que es el uso previsto en cribado virtual. La
partición por esqueleto de Bemis--Murcko impone la condición correcta. La
literatura reciente señala en la misma dirección al recomendar el desarrollo de
bancos de prueba abiertos con particiones fijas y flujos de evaluación
estandarizados, del tipo de MoleculeNet o Therapeutics Data Commons, para
permitir comparaciones justas entre estudios \cite{mao_artificial_2026}.

En segundo lugar, el dominio de aplicabilidad. El tercer principio OECD para
modelos QSAR exige declararlo explícitamente: las predicciones fuera del rango
estructural y de respuesta del conjunto de entrenamiento son extrapolaciones y
tienen menor probabilidad de ser confiables \cite{czub_automlbased_2022}. Un modelo de
cribado sin dominio de aplicabilidad declarado produce números para cualquier
entrada, incluidas aquellas sobre las que no tiene fundamento alguno.

En tercer lugar, la calidad de la evidencia de entrada. Combinar datos de
múltiples fuentes introduce el riesgo de registros contradictorios, que deben
revisarse y resolverse o eliminarse antes del modelado \cite{faramarzi_development_2022}.
Dado que ChEMBL, BindingDB y PDSP se solapan, una política de deduplicación y de
agregación de mediciones discordantes documentada de forma explícita es
condición previa a cualquier resultado interpretable.

\subsection{Aportación esperada y viabilidad}
\label{sec:just-aportacion}

La aportación de este trabajo no consiste en proponer una arquitectura nueva,
sino en construir un \textbf{procedimiento reproducible y auditable} que integre
las tres predicciones relevantes para un fármaco del SNC dirigido al receptor
5-HT\textsubscript{1A}, evaluado bajo un protocolo que refleja su uso previsto y
con su dominio de aplicabilidad declarado. En concreto, se espera entregar:

\begin{itemize}
  \item Un conjunto de datos curado y trazable de ligandos del receptor
        5-HT\textsubscript{1A}, con la política de estandarización,
        deduplicación y agregación documentada.
  \item Un protocolo de evaluación fijo y reutilizable, con particiones
        persistidas y semillas ancladas.
  \item Modelos de referencia y modelos de grafos comparados bajo idénticas
        condiciones, con sus métricas reportadas de forma completa.
  \item Un repositorio público que permita reproducir cada resultado desde los
        datos crudos.
\end{itemize}

El proyecto es viable con los recursos disponibles: todas las fuentes de datos
son públicas, todo el instrumental —RDKit, scikit-learn, PyTorch Geometric—
es de código abierto, y el cómputo requerido para conjuntos de este tamaño es
accesible en equipo de escritorio. El riesgo principal no es computacional sino
metodológico, y se concentra en la curación de los datos y en la honestidad del
protocolo de evaluación; ambos puntos son, por eso mismo, el foco del Trabajo
Terminal 1.

\todo{Si el sínodo lo requiere, añadir cifras de prevalencia nacional (ENSANUT o
INEGI) de trastornos de ansiedad y depresión en México, con cita verificada.}

